\documentclass[11pt, a4paper, sans]{article}
\usepackage[utf8]{inputenc}
\usepackage[french]{babel}
\usepackage{ragged2e}
\usepackage[parfill]{parskip}
\renewcommand{\familydefault}{\sfdefault}

\begin{document}
\thispagestyle{empty}
{\raggedright
\textbf{Joan Ko}\\3 Place Marcel Bouilloux-Lafont \\ Appartement 95\\31400 Toulouse}

{\raggedleft
\textbf{DGESIP A SCNP \\ 6 Allée Emile Monso  \\ 31405 Toulouse\\}


\vspace{5mm}
{\raggedleft
À Toulouse, le 01/10/2026}


\vspace{5mm}
\raggedright
\textbf{Objet : Candidature au poste de Data Scientist Réf. MENJ-30-2026-47058\\}
\vspace{5mm}
\noindent{

Madame, Monsieur,


Ingénieur en mathématiques appliquées diplômé de l'INSA Toulouse, je vous adresse ma candidature au poste de Data Scientist. Durant les premières années de mon parcours professionnel, j'ai pu mettre à profit mes compétences en statistiques, modélisation et analyse de données dans plusieurs domaines tels que le traitement d'image, la biologie ou l'assurance santé. En particulier, j'ai travaillé ces deux dernières années à traduire des problématiques métier en problèmes statistiques en réalisant des études actuarielles. Au cours de mes diverses expériences, j'ai appris à manipuler et traiter des jeux de données en utilisant des outils tels que Python, R ou encore Power BI.

Vous proposez un poste visant à utiliser divers outils des sciences des données tels que le machine learning pour la résolution de problématiques liées au processus d'admission Parcoursup. En particulier, cela nécessite de traiter des données à grande échelle et de construire des modèles statistiques descriptifs et prédictifs.



Ainsi, le poste que vous proposez m'intéresse particulièrement car il me permettrait de mettre à profit mes compétences en sciences des données sur plusieurs aspects. En premier lieu, je suis particulièrement intéressé par la formulation d'une problématique métier en questions statistiques. Je suis également intéressé par le prétraitement et post-traitement des données. Il s'agit, dans un premier temps, d'identifier des données pertinentes provenant de sources différentes, de réaliser leur acquisition et l'évaluation de leur fiabilité puis, dans un second temps, de les analyser pour répondre au besoin métier défini au départ. Enfin, la construction de modèles prédictifs puis leur calibration et leur évaluation sont des étapes particulièrement stimulantes pour moi. De plus, j'ai déjà eu à mettre en œuvre ces différentes étapes, suivies de l'interprétation et de la restitution des résultats, dans le cadre de mes fonctions de chargé d'études actuarielles. Par ailleurs, le cadre d'application de ce poste est motivant pour moi compte tenu de l'enjeu d'accès à l'enseignement supérieur qu'il représente.


Je me tiens à votre disposition pour toute information complémentaire. Je vous prie d'agréer, Madame, Monsieur, l'expression de mes salutations distinguées.
\vspace*{1cm}

{\raggedleft
Joan Ko

}

\end{document}

